% SHARD-ID: CA-28-GAUSSIAN-BOSON-NUMERICAL-SUITE
% SHARD-TITLE: Numerical Verification Suite for Gaussian Boson Examples
% SHARD-SUMMARY: Specifies the Julia invariant suite for scalar Gaussian boson symbols in dimensions d=1,2,3.
% SHARD-SUMMARY: Checks finite periodic stiffness matrices against Brillouin-zone symbol values instead of relying on visual spectra.
% SHARD-SUMMARY: Records the current example systems: Klein-Gordon, anisotropic cones, and doubler symbols.
% SHARD-KEYWORDS: Gaussian boson, Julia, numerical verification, finite periodic lattice, stiffness matrix, examples

\section{Numerical Verification Suite for Gaussian Boson Examples}
\label{sec:gaussian-boson-numerical-suite}

\subsection{Status, Sources, and Dependencies}

\textbf{Status: checked Julia shard.}  This shard records the numerical
invariants now implemented for the scalar Gaussian tier.  The checks are not
continuum proofs; they are finite algebraic witnesses that the coefficient
compiler is computing the intended symbols and residuals.  The massless doubler
check is a coefficient-level rejection test, not a Fock-generator check.  The
finite periodic checks below do not assert an exact differential
\([X_a,P_b]=i\delta_{ab}\) commutator.

The current numerical surface is deliberately narrow: it checks scalar
coefficient structure, finite-periodic symbol values, finite stiffness spectra,
assigned Fourier eigenvectors, periodic minimum data, and the implemented
coefficient residuals/Hessians, including sampled rotation-Hamiltonian and
boost-boost residual data on smooth-symbol momenta.  It also checks centered
periodic momentum labels for branch-aware low-energy windows, a loud finite-grid
nonnegativity gate when requested, and test-local central
finite-difference boost-time derivative oracles.  It does not test finite boost
generators, finite rotation generators, coordinate-generator commutators, or
full generator closure.

Dependencies: CA-24--CA-27, INDEX.md, and CONVENTIONS.md (j).

% Source: literature/md/2010.11121/2010.11121.md:274--293 -- finite torus lattice, dual lattice, and Fourier transform normalization.
% Source: literature/md/2010.11121/2010.11121.md:598 -- lattice scalar Hamiltonian.
% Source: literature/md/2010.11121/2010.11121.md:623 -- lattice Klein-Gordon dispersion.
% Source: literature/md/2010.11121/2010.11121.md:1626--1628 -- massless fields require a zero-mode policy.
% Checked: src/GaussianBosons.jl:9--68 and :146--263; src/GaussianBosonNumerics.jl:51--272.
% Checked: test/runtests.jl:151--220 and :225--597.

\subsection{Why These Tests}

The numerical layer is not allowed to assert that a plot ``looks relativistic''.
Each check must compare two independently constructed objects:
\[
  \text{real-space coefficient data}
  \quad\longleftrightarrow\quad
  \text{momentum-space symbol data}.
\]
For Gaussian systems this is unusually powerful because the finite periodic
stiffness matrix is exactly diagonalized by Fourier modes.

The boost-time derivative oracle follows the same rule.  The production
boost-time helper is compared to \(\frac12\nabla_k\omega(k)^2\) obtained by
central finite differences of the separate scalar-symbol evaluator.  Those
tests use generic symmetric finite-range coefficient sets in \(d=1,2,3\), with
off-axis terms in the higher-dimensional cases and sample momenta away from
high-symmetry points.

The rotation and boost-boost residual tests are also smooth-symbol checks.  They
sample coefficient formulas at named momenta and compare residual sizes or
tensor entries; they do not introduce \(X_a=i\partial_{k_a}\) as a finite
periodic matrix.

\subsection{Numerical Tolerance Policy}

CONVENTIONS.md (j) fixes the tolerance names used by this shard.  With the
\(\texttt{GAUSSIAN\_}\) prefix and \(\texttt{\_ATOL}/\texttt{\_RTOL}\) suffix
understood, the default pairs are
\[
\begin{array}{ll}
\texttt{SYMBOL\_IMAG} & (10^{-10},10^{-10}),\\
\texttt{SYMBOL\_VALUE} & (10^{-10},10^{-10}),\\
\texttt{EIGENVALUE} & (10^{-10},10^{-10}),\\
\texttt{MINIMUM\_COUNT} & (10^{-10},10^{-10}),\\
\texttt{SMALL\_SPACING\_RESIDUAL} & (5\cdot10^{-4},10^{-10}).
\end{array}
\]
The symbol-imaginary validator compares the complex symbol with its real part
after the dimension and spacing preconditions are checked.  The small-spacing
pair is a named finite-sample tolerance, not a continuum error bound.
The finite-difference derivative oracle is test-local and uses
\((2\cdot10^{-8},2\cdot10^{-8})\); it is not part of the production tolerance
surface.

\subsection{Structural Coefficient Validation}

The scalar Gaussian helpers now fail closed by default on coefficient data
outside CONVENTIONS.md (j).  The validator requires nonempty data, integer
offsets, a common dimension \(d\in\{1,2,3\}\), real coefficients, and aggregate
paired equality \(V_r=V_{-r}\) after duplicate offsets are summed.  Negative
tests cover a missing partner, unequal paired aggregates, complex coefficients,
dimension mismatch, and non-integer offsets.

\subsection{Periodic Stiffness Matrix}

For finite periodic sizes \(L_1,\ldots,L_d\), define the stiffness matrix
\[
  V_{xy}
  =
  \sum_r V_r\,\mathbf 1_{y=x+r\;{\rm mod}\;L}.
\]
For the Fourier vector
\[
  \varphi_k(x)=|\Lambda|^{-1/2}e^{i\epsilon k\cdot x},
\]
one has the exact identity
\[
  (V\varphi_k)(x)
  =
  \left(\sum_rV_re^{i\epsilon k\cdot r}\right)\varphi_k(x)
  =
  \omega_\epsilon(k)^2\varphi_k(x).
\]
Therefore each dual label \(n\) gives a tested eigenvector with its assigned
symbol value, and the sorted eigenvalues of the finite stiffness matrix must
equal the sorted symbol values on the periodic Brillouin grid.  These finite
Fourier invariants are checked in Julia for \(d=1,2,3\), including mixed
odd/even sizes with every \(L_a>2R\) and an off-axis symmetric pair
\((1,1),(-1,-1)\).  They are not finite periodic first-moment or
momentum-coordinate commutator tests; those require a branch, sawtooth,
Fourier-interpolation, or other coordinate convention before they can become
numerical claims.

The eigenvector suite also includes a one-sided shift sentinel to pin the
implementation direction \(y=x+r\).  That sentinel is not a scalar Gaussian
Hamiltonian claim.  The strict default rejects it, and the test reaches this
low-level matrix path only through an explicit non-validating keyword.

The eigenvalue oracle deliberately keeps the uncentered grid helper, because
the tested symbol spectra are branch-insensitive.  The separate centered grid
helper returns the centered integer dual label and its physical representative
\[
  k_a = \frac{2\pi n_a}{\epsilon L_a}.
\]
For even \(L=2r\), the branch is \(n_a\in\{-r,\ldots,0,\ldots,r-1\}\), matching
the centered dual lattice in the OAR source.  Thus, at \(L=4\) and
\(\epsilon=1\), the uncentered representative \(3\pi/2\) is recorded on the
centered branch as \(n=-1\) and \(k=-\pi/2\).  This is only a momentum-label
helper; it is not a periodic position-coordinate or generator construction.

\subsection{Implemented Example Systems}

\textbf{Nearest-neighbour Klein-Gordon.}  The tests check the coefficient
symbol against
\[
  m^2+\frac{2}{\epsilon^2}\sum_a(1-\cos(\epsilon k_a)),
\]
the boost-time symbol against \(\sin(\epsilon k_a)/\epsilon\), the residual
vector against \(\sin(\epsilon k_a)/\epsilon-k_a\), and the low-energy Hessian
residual against zero.  The sampled rotation and boost-boost residual helpers
are small on fixed low-energy momenta as \(\epsilon\) is made small.  The
boost-time helper is also checked on non-Klein-Gordon finite-range symbols by
the independent finite-difference oracle above.

\textbf{Anisotropic cones.}  The tests build coefficients with speeds
\((c_1,c_2,c_3)\).  They verify that the Hessian is
\[
  2\,{\rm diag}(c_1^2,c_2^2,c_3^2)
\]
and that the Lorentz residual is nonzero when the target speed is \(1\).  A
separate sampled residual test checks that anisotropy also gives nonzero
rotation-Hamiltonian and boost-boost residual data at a generic smooth momentum.

\textbf{Hessian-passing unstable scalar.}  The tests build the one-dimensional
coefficient witness
\[
  V_0=0,\qquad V_{\pm1}=1,\qquad V_{\pm2}=-\frac12.
\]
Its speed-one Hessian residual vanishes, but finite periodic sampling at
\(L=4\) has minimum \(-3\).  The minimum-data helper records one minimum at
location \(3\), uncentered label \(2\), centered label \(-2\), uncentered
momentum \(\pi\), and centered momentum \(-\pi\).  Requesting a nonnegative
finite-grid symbol raises an error, so this example cannot pass as a stable
scalar symbol.

\textbf{Doubler symbol.}  The tests build
\[
  \omega_\epsilon(k)^2
  =
  m^2+\epsilon^{-2}\sum_a\sin^2(\epsilon k_a)
\]
with \(m=0\).  In \(d=1,2,3\), the speed-one Hessian residual vanishes on the
chosen patch, but the \(L_a=4\) periodic grid has \(2^d\) zero minima with
centered label coordinates in \(\{0,-2\}\).  This catches the case where the
Hessian at one minimum is correct but the compiler would need a multi-species,
sector-selection, or zero-mode policy output.  It is deliberately outside the
positive-dispersion generator statements.  The J6 symbol-residual test also
samples a doubler momentum where the rotation residual is globally nonzero and
an extra zero where \(R_a\) is nonzero relative to the selected \(k=0\) branch.

\subsection{Current Julia Surface}

The checked helper layer is:
\[
\begin{array}{ll}
\begin{gathered}\texttt{kg\_nearest\_neighbor}\\\texttt{\_coefficients}\end{gathered}
  & \text{standard scalar KG coefficients},\\
\begin{gathered}\texttt{anisotropic\_kg}\\\texttt{\_nearest\_neighbor}\\\texttt{\_coefficients}\end{gathered}
  & \text{anisotropic coefficient witness},\\
\begin{gathered}\texttt{doubler\_quadratic}\\\texttt{\_coefficients}\end{gathered}
  & \text{coefficient-level doubler witness},\\
\begin{gathered}\texttt{validate\_scalar}\\\texttt{\_coefficients}\end{gathered}
  & \text{structural scalar-coefficient validator},\\
\texttt{scalar\_quadratic\_symbol}
  & \sum_rV_re^{i\epsilon k\cdot r},\\
\begin{gathered}\texttt{boost\_time\_symbol}\\\texttt{\_from\_coefficients}\end{gathered}
  & \frac12\nabla\omega^2,\\
\begin{gathered}\texttt{boost\_time}\\\texttt{\_residual\_from}\\\texttt{\_coefficients}\end{gathered}
  & \frac12\nabla\omega^2-c^2k,\\
\begin{gathered}\texttt{rotation\_hamiltonian}\\\texttt{\_residual\_from}\\\texttt{\_coefficients}\end{gathered}
  & k_b\partial_a\omega^2-k_a\partial_b\omega^2,\\
\begin{gathered}\texttt{boost\_boost\_residual}\\\texttt{\_coefficients\_from}\\\texttt{\_coefficients}\end{gathered}
  & \text{coefficients of }R_bX_a-R_aX_b,\\
\begin{gathered}\texttt{nearest\_neighbor}\\\texttt{\_integrated\_energy}\\\texttt{\_current\_symbol}\end{gathered}
  & \text{1D integrated A4 current symbol},\\
\begin{gathered}\texttt{low\_energy\_hessian}\\\texttt{\_from\_coefficients}\end{gathered}
  & \partial_a\partial_b\omega^2(0),\\
\texttt{lorentz\_hessian\_residual}
  & \frac12{\rm Hess}(\omega^2)(0)-c^2I,\\
\texttt{periodic\_stiffness\_matrix}
  & \text{finite periodic real-space matrix},\\
\texttt{periodic\_momentum\_grid}
  & \text{uncentered spectrum grid},\\
\begin{gathered}\texttt{centered\_periodic}\\\texttt{\_momentum\_grid}\end{gathered}
  & \text{centered labels and momenta},\\
\texttt{periodic\_fourier\_vector}
  & \text{Fourier vector for a dual label},\\
\texttt{periodic\_symbol\_values}
  & \text{finite Brillouin-grid eigenvalue oracle},\\
\begin{gathered}\texttt{symbol}\\\texttt{\_minimum\_data}\end{gathered}
  & \text{finite-grid minimum value, labels, momenta, and count},\\
\begin{gathered}\texttt{count\_periodic}\\\texttt{\_symbol\_minima}\end{gathered}
  & \text{finite-grid minima counter}.
\end{array}
\]
This list is a coverage boundary, not a generator-algebra interface: no helper
currently constructs a finite coordinate generator, a finite boost, a finite
rotation, or their commutator table.

\subsection{Next Numerical Milestones}

Proposed next suite additions, none claimed as achieved here, are:
\[
\begin{array}{ll}
1.&\text{systematic sampled residual scans on shrinking low-energy windows},\\
2.&\text{comparison of open-boundary first moments with periodic symbol tests}.
\end{array}
\]
Any finite periodic coordinate-generator version of the first-moment comparison
requires a branch or Fourier-interpolation convention before it can be promoted
beyond a local experiment.
